\section{讨论}
\label{sec:discussion}

\subsection{本文主张的边界}

在讨论本文的意义之前，需要先把主张的范围划定清楚，因为 DCC-KV 涉及的三个环节来源不同。

\textbf{不属于本文贡献的部分。}将 KV 缓存压缩为可参与后续注意力的紧凑表示，并同时匹配
注意力质量与输出，这一思想来自 MIT 团队的 Attention Matching。本文 4.3 节的 Key 选择、
质量偏置拟合与 Value 回归三步，在其思想框架内实现，\textbf{本文不主张该压缩机制为原创}。
同样，块级注意力分解、在线 Softmax 的 $(m,l,o)$ 状态表示与归并算子，均是既有工作
（FlashAttention、Online Normalizer Calculation）的组成要素。本文把后者从"单设备内分块的
数值技巧"重新理解为"跨设备异步消息的顺序无关归并机制"，但算子本身不是新的。

\textbf{属于本文贡献的部分。}其一是目的端条件化的通信抽象：对每条有向边 $(s,r)$ 单独构造
$C^{s\to r}$，使同一源块对不同目的端的紧凑表示可以不同，且显式校准块级注意力质量以维持
跨块在全局 Softmax 中的竞争关系。其二是面向变长紧凑 KV 的异步 All-to-Allv 执行机制：
环形错位调度、独立通信流与计算流、以逐块到达事件建立依赖，实现通信与计算的重叠。
其三是在此设定下给出的通信/计算复杂度分析与误差分解。这三项中，第二项是唯一在既有文献中
没有直接对应物的系统设计。

这一划分意味着本文的可证伪点也是明确的：若 E5/A2 显示逐边条件化相对共享压缩没有带来
可测的精度收益，则第一项贡献失去意义；若 E5/A5 显示异步流水相对同步实现没有带来延迟收益，
则第二项贡献失去意义。二者不相互担保。

\subsection{为什么需要目的端条件化}

一个自然的替代方案是让所有目的端共享同一份压缩 KV。这种做法更简单：源设备只需构造一次，
通信量也更低（不需要按边区分）。DCC-KV 放弃了这一优惠，换取的是每边独立的表示。

这一交换是否划算，取决于一个可测量的事实：同一源块对不同目的端的注意力分布究竟差异多大。
若差异很小，逐边构造的额外开销就是浪费；若差异显著，共享压缩就会系统性地保留对某些目的端
无用、而对另一些目的端关键的条目。H1（KL $> 0.5$）正是为这一事实设定的判据，而 E3 是唯一
能回答它的实验。

值得指出的是，这一问题的答案与工作负载强相关。在因果分块设定下，靠后的目的端与靠前的
目的端所关注的远程内容本就不同；在多跳推理任务中，不同位置需要的证据片段也不同。相反，
在需要全局平均信息的摘要类任务中，各目的端的注意力分布可能高度相似，此时条件化的收益会
明显下降。因此本文不主张目的端条件化在所有任务上都有收益，而是主张其收益的上界由 Query
分布差异决定，且这一差异可测。

\subsection{异步收益的条件与代价}
\label{sec:discussion-tension}

\S\ref{sec:analysis-async} 给出的加速上界 $1 + \min(T_{\text{comm}}, T_{\text{comp}}) /
\max(T_{\text{comm}}, T_{\text{comp}})$ 说明了一件事：异步不是无条件有益的。当通信显著慢于
计算时，加速上界趋近 $1$；当计算显著慢于通信时同样趋近 $1$。只有在两者相当时，上界才
接近 $2\times$。

这一结论对本文的叙事有一个不受欢迎但必须承认的推论。\textbf{压缩本身会减小
$T_{\text{comm}}$，因而可能把系统推离"通信与计算相当"的区间，反而降低异步的边际收益。}
也就是说，压缩与异步这两个设计之间存在张力：压缩越激进，异步能够重叠掉的时间越少。
DCC-KV 的真实收益上限来自"压缩后的通信 + 重叠"这一组合，而非二者各自收益的简单相加。
在设计实验时若只单独报告压缩比与异步加速比，会掩盖这一交互。第~\ref{sec:experiment} 节的消融 A4 正是为测定该交互而设：它同时扫预算与同步模式，并据此判定 A2 与 A5 的边缘结果能否分开引用。

异步还引入若干工程代价。变小长度消息使通信模式不规则，标准 All-to-All 需要预知接收长度，
因此须增加一次元数据交换；乱序到达要求归并算子满足交换律与结合律，这在数学上由式~\eqref{eq:comm-assoc} 保证，
但在有限精度下只近似成立（E8 的目标）；独立的通信流与计算流占用额外资源，在计算侧本已饱和
的场景下可能相互争抢。这些代价不体现在复杂度表达式中，但会体现在实测数据里。

\subsection{误差界的意义与局限}

式~\eqref{eq:err-bound} 的价值在于把端到端误差拆成质量误差与输出误差两项，并显示二者均以
全局 Softmax 分母 $S(q)$ 归一化。这带来一个有实践意义的推论：远程块占比越大、误差被吸收得
越好，因此长上下文场景比短上下文更适合该方法。这也解释了为什么 E7 预期短上下文是负结果区间。

但这一误差界有三点局限。第一，它是\emph{条件性}的：推导假设
$\sum_s |\delta M_s| < S/2$。E13 已实测该条件的违反率，它不是一个小概率事件：
$\beta$ 全量时最高 $16.7\%$、不加 $\beta$ 时最高 $100\%$，
且沿预算单调下降（$B=8$ 时 $91.0\%$ 降到 $B\ge128$ 的 $0$）。
故「界成立」与「界有保证」是两件事：界可以在前置条件不满足的格子上依然成立，
但那属于经验观测而非推导结论，二者必须分开陈述。
第二，它的紧致程度\textbf{已在合成数据上被测量}：E13 得到
$\mathrm{RHS}/\mathrm{LHS}$ 中位 $2.10$、$95$ 分位 $3.51$、最小 $1.00$
（$5760$ 个 Query 中零个反例），故它是一个量级可用的界而非装饰；
但该比值由合成分布决定，与真实模型上是否同阶仍未经测量——
$\epsilon_{\text{mass}}$ 与 $\epsilon_{\text{out}}$ 的机制级经验分布
（E2 的完整曲线，以及 E2b / E10 / E11 的网格）同样只到这一步。第三，它把代表 Query 的覆盖误差隐含在
$\epsilon_{\text{mass}}$ 与 $\epsilon_{\text{out}}$ 之中，因此没有给出 $M$ 与误差之间
显式的关系。该关系目前只由实验回答：E9 在 $M$--$B$ 网格上分离出「$B$ 决定可达上界、
$M$ 决定能否触及上界」的职责划分，但这是合成数据上的结论，且 $M$ 的取值随任务与预算
移动，尚不存在与任务无关的理论公式。

\subsection{「质量相近」的判定}
\label{sec:quality-comparable}

\S\ref{sec:exp-gpu} 的 H2 要求「DCC-KV 相对共享压缩的质量提升 $\ge 1.5$ 个百分点」
\emph{且}「在质量相近前提下 prefill 加速 $\ge 1.10\times$」。前半句的比较对象是明确的；
后半句的前提「质量相近」却只是一个词：与谁相近、由谁判定、容差多大，规划文档都没有规定。
本节给出一个可操作的定义。之所以要把它当作一个独立问题，是因为\textbf{它不是一个措辞问题，
而是一个会改变结论的问题}——容差取宽一点，任何实现都能宣称自己「质量相近」。

\paragraph{参照物取精确注意力，而不是共享压缩基线。}最省事的读法是让「质量相近」与前半句
同源，即要求 DCC-KV 与共享压缩基线在质量上相近。这一读法\textbf{自我矛盾}：若两者的质量差
被限制在 $\delta$ 以内，同时又要求二者相差至少 $1.5$ 个百分点，则只有当
$\delta > 1.5$ 个百分点时才可能同时成立；而那时「提升 $1.5$ 个百分点」已经落在「相近」的
范围之内，前半句不再表达任何主张。因此参照物必须\emph{换一个}，取\textbf{精确注意力}
（未经压缩的全局注意力）。换过之后 H2 的三段含义互不重叠、可以同时成立：压缩后的 DCC-KV
与精确注意力\textbf{不劣}（压缩基本无损），DCC-KV 相对共享压缩\textbf{更高}（逐边条件化
确有收益），且 prefill \textbf{更快}（系统设计确有收益）。

\paragraph{判定主体是协议与检验程序，不是作者。}「相近」必须由\textbf{预先固定的留出评测
协议} $\mathcal{B}$ 与一个\textbf{事先声明的检验程序}共同判定，不由作者事后目测，也不按
结果挑选指标。在 $\mathcal{B}$ 下取得 DCC-KV 与精确注意力的任务质量差
$\Delta = Q_{\text{DCC}} - Q_{\text{dense}}$ 及其 $95\%$ 置信区间，判据为
$\mathrm{CI}_{\text{low}}(\Delta) > -\delta$，即\textbf{单侧非劣检验}：压缩允许有代价，
但代价不得超过 $\delta$。采用区间而非点估计，是因为点估计为 $0$ 可能只是样本噪声恰好抵消；
用区间才能把「测不出差异」与「确实无差异」分开。这里取单侧而非双侧等价，理由是该前提的
\emph{功能}是排除「以质量换速度」这一混淆——只要 DCC-KV 没有显著更差，加速比的比较就是
公平的；若它同时显著更优（$\mathrm{CI}_{\text{low}}(\Delta) > 0$），前提以更强的形式成立，
但此时必须在结果中\textbf{显式披露方向}，不得笼统写作「相近」。

\paragraph{容差 $\delta$ 被一个下界与两个上界夹住。}下界来自统计分辨率：$\delta$ 不得小于
\textbf{同一配置重复 run 的噪声底线}（可用 run 间质量差的 p95$-$p50 或 bootstrap 区间
半宽估计）。若 $\delta$ 小于噪声底线，「相近」将永远无法被建立——测不出差异与确实无差异
不可区分，判定退化为对噪声的断言。上界之一是先验可定的：\textbf{非劣边界不得大于所要检测
的效应量}，故必须有 $\delta < 1.5$ 个百分点，否则「相近」宽到把前半句所声称的提升一并
吞没。上界之二是需要数据的：$\delta$ 还必须小于\textbf{共享压缩相对精确注意力的退化量}
$\Delta_{\text{bad}} = Q_{\text{dense}} - Q_{\text{share}}$，否则共享压缩本身也会落入
「相近」的判定范围，该前提就失去了把两者区分开的能力。于是
\begin{equation}
\label{eq:delta-window}
\text{噪声底线} \;\le\; \delta \;<\; \min\{\,1.5,\ \Delta_{\text{bad}}\,\}\ \text{个百分点}.
\end{equation}
在区间内的取值依据是保守性：先保证判定不是噪声伪影（下界），再不让它失去分辨力（上界）。
实现层面，这一区间被写成一条\textbf{硬断言}——等价检验函数在 $\delta$ 越出
\eqref{eq:delta-window} 时\emph{拒绝执行}，而不是返回「不相近」。这个区别是必要的：容差
取值不当意味着\textbf{本次实验不具备判定该前提的分辨率}，与「质量确实不相近」是两个不同的
结论，前者必须报告为分辨率不足，不能被静默读成后者。

\paragraph{两个前置条件必须与加速比测量绑定。}其一，\textbf{配对性}：质量侧与性能侧必须取自
\textbf{同一组 run}（同 batch、同提示、同上下文长度、同种子）。否则「不劣」与「更快」可能
来自不同条件，二者的组合没有意义。其二，\textbf{跨上下文长度的聚合方式}必须先定：E6 在
$4$ 个上下文长度上各产出一组质量差，是取全部长度的联合区间，还是逐长度判定后取交集，会
得到不同结论。本稿采用保守规则——\textbf{四个长度上分别做非劣检验，全部通过才算前提成立}，
任一长度不通过即单独报告该长度，不得用其余长度的通过把它平均掉。

\paragraph{失效模式。}若实测噪声底线达到或超过 $1.5$ 个百分点（在长上下文任务上并非不可能），
\eqref{eq:delta-window} 为空区间，H2 在现有测量精度下\textbf{不可检验}——此时正确的处理是
报告「分辨率不足、无法判定 H2」，而不是判定 H2 不成立。若
$\mathrm{CI}_{\text{low}}(\Delta) \in (-\delta, 0]$，则前提成立而方向未定于「更优」，
此时应如实写作「与精确注意力的差异未超过 $\delta$」，不得写作「与精确注意力相当」——
后者会隐含一个并未被检验的等号。

上述规则已落到代码：\texttt{hypotheses.py} 的 \texttt{h2\_pass\_}\allowbreak\texttt{across\_}\allowbreak\texttt{lengths}
把逐长度判定聚成单一结论，且\textbf{「分辨率不足」与「未达标」分开记账}——
某个长度的前提无法判定时记入 \texttt{unresolved} 而不计入 \texttt{failed}，
整体结果为「不可判定」；\texttt{quality\_comparable\_non\_inferior} 在
$\delta$ 越出可行区间时\emph{抛错而非返回 False}，使上面那条区别成为可执行的。
E6 的结果结构已接上该入口，其两个参数以\emph{无默认值}的命令行选项暴露：
有默认值等于假称前提永远成立。

\subsection{适用场景}

综合以上分析，DCC-KV 的适用条件可以陈述为四项同时成立：
（i）上下文足够长，使 $B/L$ 足够小，压缩带来的通信节省超过其构造与摘要开销；
（ii）跨设备带宽受限，使 $T_{\text{comm}}$ 与 $T_{\text{comp}}$ 处于可比区间，异步有发挥空间；
（iii）源块对不同目的端的注意力分布差异显著，使条件化有实际收益；
（iv）任务对 token 级精确定位的敏感度不高，使压缩误差可被容忍。

四条中任一条不成立，DCC-KV 相对精确序列并行或共享压缩的收益即会消失。这一条件集合比
"面向长上下文协同推理"的泛化表述更窄，但它更可能经得起复现。

反过来说，最不利的场景是可预期的：短上下文、高带宽互联（如单机 NVLink 全连接）、
同构目的端 Query 分布、强检索任务。在这些场景下，精确 Ring Attention 或简单共享压缩
是更合理的选择。

\section{结论}
\label{sec:conclusion}

本文提出了 DCC-KV，一种面向超长文本协同推理的目的端条件化紧凑 KV 通信框架。其基本观察是：
在序列并行设定下，源设备向不同目的设备发送的信息不应相同，因为目的端的 Query 分布不同，
而源块中各条目对不同 Query 分布的重要性也不同。据此，DCC-KV 对每条有向通信边 $(s,r)$
构造专属的紧凑 KV $C^{s\to r} = (C_k^{s\to r}, \beta^{s\to r}, C_v^{s\to r})$，
在压缩的同时显式校准块级注意力质量，并通过满足交换律与结合律的在线 Softmax 状态归并，
使已到达的远程块可以被立即计算，从而实现异步通信与计算的重叠。

本文的主要工作包括四项。第一，给出协同推理场景下的目的端条件化 KV 通信抽象，使同一源块
对于不同目的端产生不同的紧凑表示。第二，设计面向变长紧凑 KV 的异步 All-to-Allv 执行机制，
通过环形错位调度与独立的通信/计算流实现重叠，并以顺序无关的状态归并消除对消息到达顺序的
依赖。第三，给出全局注意力的块级分解，并据此将端到端误差分解为质量误差与输出误差两项，
导出以全局 Softmax 分母归一化的误差上界，说明误差不随设备数发散。第四，给出完整的实验协议，
覆盖长文本检索、问答、多跳推理与困惑度任务，并显式设定负结果报告机制。

在实现层面，紧凑 KV 构造流程、在线 Softmax 归并算子、变长消息传递与同步版分布式注意力
的数值等价性均已有可在纯 CPU 环境复现的验证，具体包括：归并算子在 $\mathrm{FP64}$ 下满足
顺序无关性，与单次密集计算的差异低于 $10^{-6}$；变长 All-to-All 在两进程 gloo 环境下
逐元素正确；\emph{不压缩且不加偏置}的同步版 DCC-KV 相对精确注意力的最大绝对偏差为
$2.5\times10^{-7}$（$\mathrm{FP32}$，$L=256$），说明通信与归并路径本身数值干净，
而压缩设定下的偏差来自紧凑构造的近似。

需要明确说明的是，本文尚未报告端到端任务指标与通信性能。异步通信流水与多卡可扩展性依赖
GPU 环境，尚未执行；关于任务质量、通信收益以及与基线的优劣对比，有待第~\ref{sec:experiment}
节所述实验完成后报告。此外，紧凑 KV 的构造机制沿用 Attention Matching 的 Key 选择、
质量偏置拟合与 Value 回归思想，本文不主张该机制为原创；本文的原创贡献集中于目的端条件化的
通信抽象与变长消息的异步执行机制。

后续工作有三个方向。其一，把代表 Query 的选择从"覆盖 Query 分布"推广到"覆盖对压缩误差
最敏感的方向"，即由误差反馈驱动摘要选择，而非由 Query 几何结构驱动。其二，把逐边独立构造成
本从 $|\mathcal{E}_s|$ 倍降低，例如对 Query 分布相近的目的端做聚类后共享部分构造中间量。
其三，把异步执行机制推广到流水线并行与专家并行的混合场景，检验顺序无关归并在更复杂通信
拓扑下的适用性。

\section*{可复现性说明}

本工作的代码与实验配置见随论文发布的代码仓库；本文对应的 commit hash 与逐次变更
记录见仓库的 \texttt{docs/}\allowbreak\texttt{commit\_}\allowbreak\texttt{log.md}。
依赖版本的冻结清单（\texttt{requirements.lock}）随代码发布提供；每个 run 的完整元数据由
\texttt{ExperimentMetadata} 记录，包括 git commit、PyTorch/CUDA/NCCL 版本、
模型名称、上下文长度、GPU 数量、随机种子、压缩预算比、同步/异步模式、
代表 Query 数与投影维度、两个正则系数，以及 RoPE 外推的使用与披露标志。
所有结果由 \texttt{RunResult} 聚合，报告 median 与 p5/p95 及 bootstrap 95\% 置信区间。
纯 CPU 部分的验证可由 \texttt{pytest tests/} 直接复现；GPU 部分的最低硬件配置见
\texttt{docs/reproducibility.md} 第 4 节。若在实验中使用了位置编码外推技术，
其使用情况会在对应结果中一并披露。
